% !TeX root = ../papers/nc-p3.tex
\section{第一長方形の特別成分の無印商と共通の四次障害}
\label{sec:first-rectangle-orbits}
\begin{theorem}\label{thm:first-rectangle-orbits}
76型の $b=c=1,t=\pm1$ において、各符号で台85の正則点の添字保持同型類の商は
\[
 \{(z,h)\in\mathbb A^2:z\ne0,\ h\ne(z-1)^2,(z+1)^2\}
\]
である。台82では各符号で $d\in\mathbb A^1\setminus\{0,1\}$ が完全不変量である。
これらの四層の各点の最小周期は2で、両符号の間や両台の間に同型はない。
76型の水平主成分の指定始点の一般無印商は有理三次元で、正規化座標 $(t,b,x)$ は一般同型類を一意に与える。
\end{theorem}
\begin{proof}
台85では $B=C=1$, $x=A$, $d=D$ として $xd(x-1)(d-1)(x-d)\ne0$ となる。
台82では $x=1,d\ne0,1$ である。両方の全点の正則性は定理\ref{thm:parallel-support-fibers}で示した。

まず四文字枠の復元を行う。特別成分で $x+d\ne1$ なら位置1の四線を二次モーメントと挿入付き縮約で分離し、
次に位置3を、最後に位置0,2を復元する。台85の二変数環で各準備段階の差のイデアルは
位置1の文字対 $(0,3),(1,2)$ のみ $(x+d-1)$ で、他はすべて単位イデアルとなる。
台82の一変数環では全て単位イデアルである。

$x+d=1$ のとき、各位置 $p$ に対して異なる位置 $q$ を三通り取り、$p,q$ で切断した
$W_{pq}G_{pq}^t$ の平方の部分traceを取る。$G$ は四次交代収縮の勾配である。
各対角列は $i\leftrightarrow i\mathbin\oplus e_p$, $e=(2,3,2,3)$ で一定で、文字0と1での
三差の分子の最大公約数は定数である。これで各位置の二平面 $E_{p,j}=\langle j,j\oplus e_p\rangle$,
$j=0,1$ を内在的に区別できる。これらの平面への射影で四重テンソルを切り、位置02と13で平坦化する。
$j$ の四つ組が $(0,0,1,0)$ と $(1,1,0,1)$ のブロックは、それぞれ行と列の非零な二座標が
両添字とも異なる位置にあり、その二行二列が可逆である。従って像および双対像の射影直線と
階数一のSegre曲面の交点はそれぞれ座標の二点である。全ての平面の両文字を得る。
ブロックを特定するのに文字名の事前の選択は要らず、復元した平面の全組を取ればよい。
他のブロックも同じ文字対にのみ支持され、階数二のものはいずれも同じ座標線を与える。
同じ論証は水平主成分の $(t,b,x)=(2,3,7)$ でも通る。後述の平行移動対称性により二平面内の
重複固有値は水平主成分上で恒等的なので、この計算は一般点にも適用できる。
\path{computations/verify_first_rectangle_frames.py} が両符号のイデアルと二次元切断の行列式を検算する。

従って全ての同型は単項である。特別成分の台85で指定始点を保つ台置換は二つあり、
係数では恒等と $(x,d)\mapsto(1-d,1-x)$ となる。一位置移動は $(x,d)\mapsto(d,x)$ である。
$z=x-d$, $s=x+d-1$ を使うと指定始点の作用は $(z,s)\mapsto(z,-s)$ であり、$h=s^2$ で
上の商を得る。一位置移動で $z\mapsto-z$ となるので全点の周期は2である。
台82では偶数移動ごとに一つの台置換があり $d$ を保ち、奇数移動の台置換はない。
どちらも $t$ の符号を変える指定始点の台置換はない。

水平主成分では $x=A$, $B=C=1$, $D=c(1-x)/m$, $m=t^2=b+c-bc$ と置く。
台85の指定始点の二置換はともに正規化係数を変えない。一位置移動とその平方は
\[
 J(t,b,c,x)=(1/t,c/m,1/b,c(1-x)/m),\qquad
 J^2(t,b,c,x)=(t,m/b,m/c,1-c+cx/b)
\]
となる。$J^4=1$ である。これらは分母の許す全点で同型を実現する。
枠を復元できる開集合ではこれが全同型である。行列の右下は薄い係数と左上から一意なので、
薄い指数の規格化と左上八係数の比較で十分である。
\path{computations/verify_first_rectangle_orbits.py} が全置換を検算する。
\end{proof}

\begin{proposition}
水平一般点、特別成分の台85の $x+d=1$ 上では、窓の直線の射影安定群は $(\mathbb Z/2)^3$ である。
そのうち四位置の全てで可換な線型持上げを持つ位数4の部分群は標準の四元群一つだけである。
台85の他の特別成分点と台82の特別成分点では射影安定群は標準の四元群である。
\end{proposition}
\begin{proof}
両台とも位置別重み行列の非零Smith因子は $1^{11},2,2$ で、射影対角群 $K_0$ は位数4である。
追加の置換が係数を保つ場合、その平行移動量は各位置で $e=(2,3,2,3)$ である。
その単項持上げ $P$ は文字符号と射影的に可換であり、$P^2\in K_0$ の対角値は各 $e_p$ の対で一定となる。
$2,3$ が四元群を生成するので、全対で一定な文字は自明である。従って射影群は初等可換2群となる。
他の位数4の部分群は $K_0$ と非自明な文字 $\chi$ を共有する。対角外の元とその文字の線型交換子は
位置 $p$ で $\chi(e_p)$ であり、全位置で1にはならない。この交換子は持上げのスカラーを変えても変わらない。
\end{proof}

\begin{proposition}\label{prop:parallel-antidiagonal-obstruction}
76型・112型いずれの主成分でも $b+c=0$ の全点は標準AS正則性を失う。
特に台143における $y^2-y-1=0$ と台145における $ay=-1$ は同じ四次障害の切片である。
\end{proposition}
\begin{proof}
$c=-b$ なら $t=\pm b$ である。各符号で76型は自由な $b,x$ を、
112型は自由な $b,a,y$ を持ち、補完式の定義域は既約である。始点0で一般係数の
長さ4の四文字の経路次元を計算すると、四つの一般点でいずれも $(10,10,8,8)$ である。
その和36は標準値35より大きく、二次関係が一定階数の範囲での上半連続性により特殊点でも排除される。
\path{computations/verify_rectangle_antidiagonal_obstruction.py} が四有理関数体で計算する。
76型の上記の枠復元開集合で $J^2$ が係数を固定するには $b^2=c^2=m$ が必要で、
$c=b$ なら $b=c=1$、$c=-b$ ならこの障害を受ける。従ってこの開集合の正則点は
特別成分との交線では周期2で、その外では周期4である。
\end{proof}
